\chapter*{General Introduction}
\addcontentsline{toc}{chapter}{General Introduction}
\markboth{General Introduction}{General Introduction}

\section*{Context}
Football clubs increasingly depend on digital systems to manage operational
activities such as documents, events, staff coordination, player information,
and internal communication. In practice, these activities are often distributed
across independent tools with different data models and permission mechanisms.
This situation reduces operational coherence and complicates controlled
collaboration between administrative, technical, and medical stakeholders.

\section*{Problem Statement}
The central problem addressed in this project is the following: how to design a
unified and modular platform that centralizes football club operations while
ensuring traceability, consistent access control, and maintainable service
separation. Existing tool fragmentation generally leads to scattered data,
limited audit visibility, and heterogeneous authorization practices that are not
aligned with club roles and responsibilities.

\section*{Project Objectives}
The project objective is to design and implement \textbf{Football OS} as a
structured operational workspace based on a microservices-oriented
architecture. The
solution combines an Angular frontend with a Spring Boot backend organized
around an API Gateway, a Governance Service, a Workspace Service, and a Shared
SDK.

The SMART-oriented objectives adopted for this work are summarized below:
\begin{enumerate}
    \item Centralize operational documents and events in a unified workspace model,
    \item Secure access through authentication and role-based authorization rules,
    \item Provide player profile consultation filtered by granted permissions,
    \item Integrate online document viewing and editing workflows,
    \item Structure backend responsibilities around Gateway, Governance, Workspace, and Shared SDK modules.
\end{enumerate}

\section*{Adopted Methodology}
The work followed an iterative incremental approach organized around successive
technical layers: backend foundation, governance mechanisms, workspace feature
implementation, frontend integration, and validation. Each iteration delivered
an integrated subset of capabilities and then consolidated routing,
authorization, data handling, and user workflows before extending the scope.

\section*{Report Organization}
This report is organized into five main chapters. Chapter~1 introduces the
internship context and project positioning. Chapter~2 formalizes requirements
through actors, functional and non-functional constraints, and use-case
contracts. Chapter~3 analyzes system design decisions across architecture,
module boundaries, data modeling, and processing flows. Chapter~4 documents
implementation, integration, interfaces, and deployment choices. Chapter~5
presents testing results, validation limits, and improvement priorities. The
report ends with a general conclusion.

